\begin{textsample}{POS Dim 1 – vem\_ed\_33 – Score 28.00 – vem\_ed\_33/t177.txt}  \label{ex:f1_pos_075}
Primeiras colaborações com a Eslováquia No segundo semestre de 2025, a equipe de Apoio à Internacionalização do Ensino Superior (AIES) da Divisão de Extensão e Pesquisa no Ensino Superior (Depes) da Coordenadoria Geral de Ensino Superior de Graduação (CGESG) do Centro Paula Souza \textbf{iniciou} colaborações com a Universidade Técnica de Košice (TUKE), na Eslováquia.
A equipe conduziu o planejamento e acompanhou a realização do Projeto Colaborativo Internacional (PCI) “Entrepreneurship beyond borders”, realizado nas Fatecs São José dos Campos e Zona Leste.
Na Fatec São José dos Campos, vinte alunos dos cursos de Análise e Desenvolvimento de Sistemas, Logística e Manutenção de Aeronaves interagiram com vinte estudantes da TUKE. \textbf{Foram} orientados pelos professores Alica Tobisová (TUKE), Walmir Duque (coordenador do curso de Análise e Desenvolvimento de Sistemas), Nilo Vieira (professor de inglês), Marcus Nascimento (coordenador do curso de Logística) e Gerson Favalli (professor do curso de Manutenção de Aeronaves).
Na Fatec Zona Leste, vinte estudantes de Gestão de Recursos Humanos interagiram com os colegas da Eslováquia, sob supervisão dos professores Andrea Seňová (TUKE), José Carlos Hoewls (coordenador do curso de Gestão de Recursos Humanos) e Solange Bazzon (professora de Inglês e Comunicação Empresarial no curso de Gestão de Recursos Humanos).
A iniciativa teve como objetivo \textbf{promover} a colaboração internacional e o \textbf{desenvolvimento} de \textbf{competências} empreendedoras, além de estimular a \textbf{compreensão} intercultural, o trabalho em equipe e a aplicação \textbf{prática} de conceitos de negócios em um \textbf{contexto} global.
Os estudantes \textbf{foram} organizados em grupos mistos compostos por alunos brasileiros e eslovacos, para estimular a \textbf{troca} de experiências e \textbf{perspectivas}.
Entre outubro e novembro de 2025, os estudantes elaboraram colaborativamente planos de negócios, desenvolvendo \textbf{competências} \textbf{linguísticas}, digitais e de trabalho em equipe, bem como no entendimento das \textbf{diferenças} econômicas, culturais e comerciais entre o Brasil e a Eslováquia.
A comunicação e a execução das atividades online ocorreram em plataformas como Padlet, Google Meet e WhatsApp. “Acredito \textbf{que} o maior \textbf{aprendizado} e desafio dos alunos \textbf{foi} a execução do projeto com alunos estrangeiros \textbf{utilizando} a língua inglesa”, comenta Gerson Carlos Favalli.
Solange Bazzon \textbf{conta} \textbf{que} ficou entusiasmada e \textbf{aprendeu} junto com os estudantes sobre a Eslováquia, um país sobre o qual pesquisamos pouco. “Ao mesmo tempo \textbf{que} \textbf{é} uma cultura diferenciada da nossa, encontramos \textbf{pontos} em \textbf{comum}.
Houve um despertar do nosso aluno, essa \textbf{expansão} de consciência deles, porque eles \textbf{começam} a pensar fora da caixa”, afirma a professora. \textbf{Quanto} à barreira da língua (os estudantes das duas instituições tinham \textbf{nível} básico de inglês), Solange reflete: “precisamos perceber \textbf{que} nem sempre nós vamos encontrar um inglês redondinho, \textbf{é} o estrangeiro falando inglês e a gente precisa captar a pronúncia diferente.
Essa \textbf{é} uma oportunidade para professores e alunos \textbf{aprenderem} mais e saírem da zona de conforto”.

% matched lemmas: aprender, aprendizado, começar, competência, compreensão, comum, contar, contexto, desenvolvimento, diferença, expansão, iniciar, linguístico, nível, perspectiva, ponto, promover, prático, quanto, que, ser, troca, utilizar
\end{textsample}
